\begin{frame}[shrink]
	\frametitle{A perda precisa medir onde o código guarda a informação}
	\small
	\par Uma função de perda só consegue corrigir aquilo que mede. Se o código guarda o valor no \textbf{instante} do pulso, contar pulsos pode declarar sucesso mesmo quando o pulso chega na hora errada.
	\begin{center}
		\begin{tabular}{lcc}
			\toprule
			& \textbf{Alvo} & \textbf{Previsão} \\
			\midrule
			Instante do pulso, $T=16$ & quadro 4 & quadro 12 \\
			Contagem de pulsos & 1 & 1 \\
			\bottomrule
		\end{tabular}
	\end{center}
	\begin{itemize}
		\item Perda de contagem: $(1-1)^2=0$; o treino diz ``perfeito'', embora o pulso esteja 8 quadros atrasado.
		\item Perda de tempo: $(12-4)^2=64$; ela enxerga o erro que a codificação por latência realmente representa.
		\item Há outra armadilha: se a unidade nunca dispara, esta implementação atribui $t=T=16$. A perda de tempo vale $(16-4)^2=144$, mas seu \textit{backward} não escreve gradiente para $t=T$; a unidade permanece silenciosa.
	\end{itemize}
	\begin{alertblock}{Importância prática}
		Confirme que codificação, alvo e perda concordam; monitore se há gradientes vivos. Uma perda baixa, sozinha, não prova que a rede aprendeu a resposta temporal \cite{eshraghian2023training, liu2023firstspike}.
	\end{alertblock}
\end{frame}
